\documentclass{article}
\usepackage{amsmath,amssymb,amsthm}
\usepackage{algorithm}
\usepackage[noend]{algpseudocode}
\usepackage{enumitem}
\usepackage{hyperref}
\newtheorem{theorem}{Theorem}
\newtheorem{lemma}{Lemma}
\newtheorem{assumption}{Assumption}
\newcommand{\E}{\mathbb{E}}
\newcommand{\R}{\mathbb{R}}

\begin{document}

\section*{Federated Averaging with Local SGD (FedAvg-LS)}

\section*{Mathematical Formulation}
Let there be $K$ participating clients.  
Each client $k\in\{1,\dots,K\}$ possesses a local convex $L$‑smooth loss 
\[
f_k(w)=\E_{\xi\sim\mathcal{D}_k}\bigl[\,\ell(w;\xi)\,\bigr],\qquad 
f(w)=\frac1K\sum_{k=1}^{K}f_k(w)
\]
and can draw i.i.d. stochastic samples $\xi_{k}^{t,e}\sim\mathcal{D}_k$.

\medskip
\noindent\textbf{Hyper‑parameters}
\begin{itemize}[leftmargin=*, nosep]
  \item Learning rate (step size) $\eta>0$ (constant for the whole run).
  \item Number of communication rounds $T\in\mathbb{N}$.
  \item Number of local SGD steps per round $E\in\mathbb{N}$.
\end{itemize}

\medskip
\noindent\textbf{Iterates}
\[
\begin{aligned}
&\text{global model at round }t: && w^{t}\in\R^{d},\\
&\text{local model of client }k \text{ after }e\text{ local steps: } 
&& w_{k}^{t,e}\in\R^{d},\qquad e=0,\dots,E .
\end{aligned}
\]

\medskip
\noindent\textbf{Update rules}
\begin{enumerate}
  \item \textbf{Broadcast:} at the beginning of round $t$ the server sends $w^{t}$ to every client.
  \item \textbf{Local SGD:} each client $k$ initializes $w_{k}^{t,0}=w^{t}$ and for $e=0,\dots,E-1$ performs
  \[
  g_{k}^{t,e}\;:=\;\nabla \ell\bigl(w_{k}^{t,e};\xi_{k}^{t,e}\bigr),\qquad
  w_{k}^{t,e+1}\;=\;w_{k}^{t,e}-\eta\,g_{k}^{t,e}.   \tag{1}
  \]
  \item \textbf{Upload \& Averaging:} after $E$ steps client $k$ sends $w_{k}^{t,E}$ to the server.
  The server updates the global model by averaging
  \[
  w^{t+1}\;=\;\frac1K\sum_{k=1}^{K}w_{k}^{t,E}.   \tag{2}
  \]
\end{enumerate}
Define the \emph{client drift} at round $t$ as
\[
\delta^{t}\;:=\;\frac1K\sum_{k=1}^{K}\frac1E\sum_{e=0}^{E-1}
\bigl(\nabla f_k(w_{k}^{t,e})-\nabla f(w^{t})\bigr). \tag{3}
\]

\section*{Pseudocode}
\begin{algorithm}[H]
\caption{FedAvg-LS (Local SGD with $E$ steps)}\label{alg:fedavg}
\begin{algorithmic}[1]
\State \textbf{Input:} step size $\eta$, local steps $E$, rounds $T$, initial model $w^{0}$
\For{$t=0$ {\bfseries to} $T-1$}
    \State Server broadcasts $w^{t}$ to all clients
    \For{each client $k=1,\dots,K$ \textbf{in parallel}}
        \State $w_{k}^{t,0}\gets w^{t}$
        \For{$e=0$ {\bfseries to} $E-1$}
            \State Sample $\xi_{k}^{t,e}\sim\mathcal{D}_k$
            \State $g_{k}^{t,e}\gets\nabla\ell\bigl(w_{k}^{t,e};\xi_{k}^{t,e}\bigr)$
            \State $w_{k}^{t,e+1}\gets w_{k}^{t,e}-\eta\,g_{k}^{t,e}$ \Comment{Eq.~(1)}
        \EndFor
        \State Send $w_{k}^{t,E}$ to server
    \EndFor
    \State $w^{t+1}\gets\frac1K\sum_{k=1}^{K}w_{k}^{t,E}$ \Comment{Eq.~(2)}
\EndFor
\State \textbf{Output:} $w^{T}$
\end{algorithmic}
\end{algorithm}

\section*{Dry Run Example}
Consider the scalar quadratic $f(w)=(w-3)^{2}$ (hence $K=1$, $E=2$).  
Stochastic gradient is exact: $g^{t,e}= \nabla f(w^{t,e}) = 2(w^{t,e}-3)$.  
Take $\eta=0.1$, $w^{0}=0$.

\begin{center}
\begin{tabular}{c|c|c|c}
\textbf{Step} & $w^{t,e}$ & $g^{t,e}=2(w^{t,e}-3)$ & $f(w^{t,e})$\\\hline
$t=0,e=0$ & $0$ & $-6$ & $9$\\
$t=0,e=1$ & $0-0.1(-6)=0.6$ & $2(0.6-3)=-4.8$ & $5.76$\\
$t=0,e=2$ & $0.6-0.1(-4.8)=1.08$ & $2(1.08-3)=-3.84$ & $3.6864$\\
\multicolumn{4}{c}{Average (single client) $\Rightarrow w^{1}=1.08$}\\\hline
$t=1,e=0$ & $1.08$ & $2(1.08-3)=-3.84$ & $3.6864$\\
$t=1,e=1$ & $1.08-0.1(-3.84)=1.464$ & $2(1.464-3)=-3.072$ & $2.365$\\
$t=1,e=2$ & $1.464-0.1(-3.072)=1.7712$ & $2(1.7712-3)=-2.4576$ & $1.511$\\
\multicolumn{4}{c}{$\Rightarrow w^{2}=1.7712$}
\end{tabular}
\end{center}
After three communication rounds the global model moves progressively towards the optimum $w^{*}=3$, and the objective values decrease from $9$ to $1.51$.

\newpage
\section*{Assumptions}
\begin{assumption}[Convexity and Smoothness]\label{ass:convex}
Each local loss $f_k$ is convex and $L$‑smooth, i.e.
\[
f_k(y)\ge f_k(x)+\langle\nabla f_k(x),y-x\rangle,\qquad
\|\nabla f_k(y)-\nabla f_k(x)\|\le L\|y-x\|,\;\forall x,y\in\R^{d}.
\]
Consequently $f$ inherits the same properties.
\end{assumption}

\begin{assumption}[Unbiased Stochastic Gradients]\label{ass:unbiased}
For any client $k$, round $t$, local step $e$,
\[
\E\!\left[g_{k}^{t,e}\mid w_{k}^{t,e}\right]=\nabla f_k\!\left(w_{k}^{t,e}\right).
\]
\end{assumption}

\begin{assumption}[Bounded Variance]\label{ass:variance}
There exists $\sigma^{2}>0$ such that
\[
\E\!\left[\|g_{k}^{t,e}-\nabla f_k(w_{k}^{t,e})\|^{2}\mid w_{k}^{t,e}\right]\le\sigma^{2},
\qquad\forall k,t,e.
\]
\end{assumption}

\begin{assumption}[Bounded Drift]\label{ass:drift}
Define the drift vector $\delta^{t}$ as in (3).  
Assume a uniform bound
\[
\|\delta^{t}\|\le\Delta,\qquad\forall t,
\]
where $\Delta\ge0$ quantifies the heterogeneity among clients.
\end{assumption}

\begin{assumption}[Step Size]\label{ass:stepsize}
The learning rate satisfies
\[
0<\eta\le \frac{1}{2L\,(E+1)}.
\]
\end{assumption}

\section*{Main Convergence Theorem}
\begin{theorem}\label{thm:conv}
Under Assumptions \ref{ass:convex}–\ref{ass:stepsize}, let $w^{*}\in\arg\min_{w}f(w)$.  
For the sequence $\{w^{t}\}_{t=0}^{T}$ generated by FedAvg-LS,
\[
\frac{1}{T}\sum_{t=0}^{T-1}\E\!\left[f\!\left(w^{t}\right)-f\!\left(w^{*}\right)\right]
\;\le\;
\underbrace{\frac{\|w^{0}-w^{*}\|^{2}}{2\eta T}}_{\text{optimization term}}
\;+\;
\underbrace{ \frac{L\eta\sigma^{2}}{K}}\_{\text{stochastic term}}
\;+\;
\underbrace{L\eta\Delta^{2}E}_{\text{drift term}}.
\tag{4}
\]
Consequently, choosing $\eta = \Theta\!\bigl(\tfrac{1}{\sqrt{T}}\bigr)$ yields the rate
\[
\E\!\left[f(\bar w_{T})-f(w^{*})\right]=O\!\left(\frac{1}{\sqrt{T}}+\frac{E\Delta^{2}}{\sqrt{T}}\right),
\qquad\bar w_{T}:=\frac1T\sum_{t=0}^{T-1}w^{t}.
\]
\end{theorem}

\section*{Theoretical Properties}
\begin{itemize}[leftmargin=*]
  \item \textbf{Stability w.r.t.\ step size.} The bound \eqref{eq:stepsize} guarantees that the local descent does not overshoot the curvature of each $f_k$.
  \item \textbf{Effect of client heterogeneity.} The term $L\eta\Delta^{2}E$ grows linearly with the number of local steps $E$ and the drift magnitude $\Delta$, reflecting the well‑known “client drift’’ phenomenon.
  \item \textbf{Comparison to centralized SGD.} Setting $K=1$ and $E=1$ reduces \eqref{4} to the classical SGD bound $\|w^{0}-w^{*}\|^{2}/(2\eta T)+L\eta\sigma^{2}$.
  \item \textbf{Benefit of many clients.} The stochastic variance term scales as $1/K$, showing variance reduction from parallelism.
\end{itemize}

\section*{Discussion}
The theorem shows that FedAvg‑LS attains the optimal $O(1/\sqrt{T})$ rate for convex smooth problems up to an additive drift penalty that is unavoidable when data are non‑identically distributed. By decreasing $\eta$ as $1/\sqrt{T}$, both the optimization and drift contributions decay at the same rate, while the variance term is already $O(1/T)$ after averaging over $K$ devices.

\newpage
\section*{Preliminaries (Definitions and Lemmas)}
\begin{lemma}[Smoothness inequality]\label{lem:smooth}
For any $L$‑smooth function $h$,
\[
h(y)\le h(x)+\langle\nabla h(x),y-x\rangle+\frac{L}{2}\|y-x\|^{2},\qquad\forall x,y.
\]
\end{lemma}
\begin{proof}
Standard consequence of $L$‑smoothness; see e.g. \cite{Nesterov2004}. \qedhere
\end{proof}

\begin{lemma}[Three‑point identity]\label{lem:three}
For any vectors $a,b,c$,
\[
\|a-c\|^{2}= \|a-b\|^{2}+2\langle a-b,c-b\rangle+\|c-b\|^{2}.
\]
\end{lemma}
\begin{proof}
Expand both sides and cancel terms. \qedhere
\end{proof}

\section*{Main Proof (Step‑by‑Step Derivation)}
We prove Theorem \ref{thm:conv}.  All expectations are taken over the randomness of the stochastic gradients and, when convenient, over the random choice of the client set (here all $K$ clients participate every round).

\medskip
\noindent\textbf{Notation.} Throughout we denote $w^{t}$ the global model before round $t$, and $w_{k}^{t,e}$ the local model of client $k$ after $e$ local steps in round $t$.  Define the averaged local model after $E$ steps,
\[
\bar w^{t}:=\frac1K\sum_{k=1}^{K} w_{k}^{t,E}=w^{t+1}.
\]

\medskip
\noindent\underline{[Step 1]} \textbf{Smoothness of $f$ applied to the global update.}  
Using Lemma \ref{lem:smooth} with $x=w^{t}$ and $y=w^{t+1}=\bar w^{t}$,
\[
f(w^{t+1})\le f(w^{t})+\langle\nabla f(w^{t}),\,w^{t+1}-w^{t}\rangle+\frac{L}{2}\|w^{t+1}-w^{t}\|^{2}. \tag{5}
\]

\medskip
\noindent\underline{[Step 2]} \textbf{Express $w^{t+1}-w^{t}$ via local updates.}  
From the definition of $\bar w^{t}$ and the local SGD rule (1),
\[
w^{t+1}-w^{t}= \frac1K\sum_{k=1}^{K}\bigl(w_{k}^{t,E}-w_{k}^{t,0}\bigr)
 = -\eta\frac1K\sum_{k=1}^{K}\sum_{e=0}^{E-1} g_{k}^{t,e}. \tag{6}
\]

\medskip
\noindent\underline{[Step 3]} \textbf{Plug (6) into (5).}  
\[
\begin{aligned}
f(w^{t+1}) &\le f(w^{t}) -\eta\Big\langle\nabla f(w^{t}),\frac1K\sum_{k}\sum_{e} g_{k}^{t,e}\Big\rangle
 +\frac{L\eta^{2}}{2}\Big\|\frac1K\sum_{k}\sum_{e} g_{k}^{t,e}\Big\|^{2}.
\end{aligned} \tag{7}
\]

\medskip
\noindent\underline{[Step 4]} \textbf{Take expectation conditioned on $w^{t}$.}  
Using Assumption \ref{ass:unbiased} and linearity of expectation,
\[
\E\!\bigl[g_{k}^{t,e}\mid w^{t}\bigr]=\E\!\bigl[\nabla f_k(w_{k}^{t,e})\mid w^{t}\bigr].
\]
Hence
\[
\E\!\Big[\frac1K\sum_{k}\sum_{e} g_{k}^{t,e}\,\Big|\,w^{t}\Big]
 = \sum_{e=0}^{E-1}\frac1K\sum_{k}\nabla f_k(w_{k}^{t,e})
 = E\,\nabla f(w^{t}) + E\,\delta^{t}, \tag{8}
\]
where the last equality follows from the definition of $\delta^{t}$ in (3).

\medskip
\noindent\underline{[Step 5]} \textbf{Bound the inner product term.}  
Using (8) and the identity $\langle a,b\rangle = \frac12\bigl(\|a\|^{2}+\|b\|^{2}-\|a-b\|^{2}\bigr)$,
\[
\begin{aligned}
&\E\!\Big[\langle\nabla f(w^{t}),\tfrac1K\sum_{k}\sum_{e} g_{k}^{t,e}\rangle\mid w^{t}\Big]\\
&= \langle\nabla f(w^{t}),\,E\nabla f(w^{t})+E\delta^{t}\rangle\\
&= E\|\nabla f(w^{t})\|^{2}+E\langle\nabla f(w^{t}),\delta^{t}\rangle\\
&\ge E\|\nabla f(w^{t})\|^{2}-E\|\nabla f(w^{t})\|\,\|\delta^{t}\| \\
&\ge E\|\nabla f(w^{t})\|^{2}-\frac{E}{2}\bigl(\|\nabla f(w^{t})\|^{2}+ \|\delta^{t}\|^{2}\bigr)\\
&= \frac{E}{2}\|\nabla f(w^{t})\|^{2}-\frac{E}{2}\|\delta^{t}\|^{2}.
\end{aligned} \tag{9}
\]

\medskip
\noindent\underline{[Step 6]} \textbf{Bound the quadratic term.}  
From Lemma \ref{lem:three} and Jensen's inequality,
\[
\begin{aligned}
&\E\!\Big\|\frac1K\sum_{k}\sum_{e} g_{k}^{t,e}\Big\|^{2}
 = \E\!\Big\| \sum_{e=0}^{E-1}\bigl(\frac1K\sum_{k}g_{k}^{t,e}\bigr)\Big\|^{2}\\
&\le E\sum_{e=0}^{E-1}\E\!\Big\|\frac1K\sum_{k}g_{k}^{t,e}\Big\|^{2}
 \quad\text{(Cauchy–Schwarz for sums)}\\
&=E\sum_{e=0}^{E-1}\E\!\Big\|\frac1K\sum_{k}\bigl(\nabla f_k(w_{k}^{t,e})+ \underbrace{g_{k}^{t,e}-\nabla f_k(w_{k}^{t,e})}_{\text{noise}}\bigr)\Big\|^{2}\\
&\le 2E\sum_{e=0}^{E-1}\Bigg(
      \underbrace{\E\!\Big\|\frac1K\sum_{k}\nabla f_k(w_{k}^{t,e})\Big\|^{2}}_{\text{bias}}
      +\underbrace{\E\!\Big\|\frac1K\sum_{k}\bigl(g_{k}^{t,e}-\nabla f_k(w_{k}^{t,e})\bigr)\Big\|^{2}}_{\text{variance}}
    \Bigg).
\end{aligned} \tag{10}
\]

\medskip
\noindent\underline{[Step 7]} \textbf{Bias term.}  
Using convexity of the squared norm and the definition of $\delta^{t}$,
\[
\begin{aligned}
\E\!\Big\|\frac1K\sum_{k}\nabla f_k(w_{k}^{t,e})\Big\|^{2}
&= \Big\|\frac1K\sum_{k}\bigl(\nabla f_k(w_{k}^{t,e})-\nabla f(w^{t})\bigr)+\nabla f(w^{t})\Big\|^{2}\\
&\le 2\Big\|\nabla f(w^{t})\Big\|^{2}+2\Big\|\frac1K\sum_{k}\bigl(\nabla f_k(w_{k}^{t,e})-\nabla f(w^{t})\bigr)\Big\|^{2}\\
&\le 2\|\nabla f(w^{t})\|^{2}+2\Delta^{2},
\end{aligned} \tag{11}
\]
where we used Assumption \ref{ass:drift} and the fact that the average over $e$ does not increase the bound.

\medskip
\noindent\underline{[Step 8]} \textbf{Variance term.}  
Since stochastic gradients are independent across clients and steps,
\[
\begin{aligned}
\E\!\Big\|\frac1K\sum_{k}\bigl(g_{k}^{t,e}-\nabla f_k(w_{k}^{t,e})\bigr)\Big\|^{2}
&= \frac1{K^{2}}\sum_{k}\E\!\Big\|g_{k}^{t,e}-\nabla f_k(w_{k}^{t,e})\Big\|^{2}\\
&\le \frac{\sigma^{2}}{K}.
\end{aligned} \tag{12}
\]

\medskip
\noindent\underline{[Step 9]} \textbf{Combine (10)–(12).}  
Plugging the bounds (11) and (12) into (10) yields
\[
\begin{aligned}
\E\!\Big\|\frac1K\sum_{k}\sum_{e} g_{k}^{t,e}\Big\|^{2}
&\le 2E\sum_{e=0}^{E-1}\Bigl(2\|\nabla f(w^{t})\|^{2}+2\Delta^{2}+\frac{\sigma^{2}}{K}\Bigr)\\
&= 4E^{2}\|\nabla f(w^{t})\|^{2}+4E^{2}\Delta^{2}+ \frac{2E^{2}\sigma^{2}}{K}.
\end{aligned} \tag{13}
\]

\medskip
\noindent\underline{[Step 10]} \textbf{Insert (9) and (13) into (7).}  
Taking total expectation and using $\E[\cdot\mid w^{t}]$ notation,
\[
\begin{aligned}
\E\!\bigl[f(w^{t+1})\mid w^{t}\bigr]
&\le f(w^{t})-\eta\Bigl(\frac{E}{2}\|\nabla f(w^{t})\|^{2}-\frac{E}{2}\|\delta^{t}\|^{2}\Bigr)\\
&\qquad+\frac{L\eta^{2}}{2}\Bigl(4E^{2}\|\nabla f(w^{t})\|^{2}+4E^{2}\Delta^{2}+ \frac{2E^{2}\sigma^{2}}{K}\Bigr).
\end{aligned} \tag{14}
\]

\medskip
\noindent\underline{[Step 11]} \textbf{Choose $\eta$ small enough.}  
Assumption \ref{ass:stepsize} guarantees $L\eta E\le \tfrac12$, i.e. $2L\eta E\le1$.  
Therefore,
\[
\frac{L\eta^{2}}{2}\,4E^{2}\|\nabla f(w^{t})\|^{2}
\le \frac{\eta E}{2}\|\nabla f(w^{t})\|^{2}. \tag{15}
\]

\medskip
\noindent\underline{[Step 12]} \textbf{Simplify (14) using (15).}  
\[
\begin{aligned}
\E\!\bigl[f(w^{t+1})\mid w^{t}\bigr]
&\le f(w^{t})-\frac{\eta E}{2}\|\nabla f(w^{t})\|^{2}
   +\frac{\eta E}{2}\|\delta^{t}\|^{2}\\
&\qquad+ 2L\eta^{2}E^{2}\Delta^{2}+ \frac{L\eta^{2}E^{2}\sigma^{2}}{K}.
\end{aligned} \tag{16}
\]

\medskip
\noindent\underline{[Step 13]} \textbf{Relate gradient norm to suboptimality.}  
Convexity of $f$ gives $f(w^{t})-f(w^{*})\le \langle\nabla f(w^{t}),w^{t}-w^{*}\rangle\le \|w^{t}-w^{*}\|\|\nabla f(w^{t})\|$.  
Using Young’s inequality $ab\le \frac{1}{2\eta E}\,a^{2}+\frac{\eta E}{2}\,b^{2}$ with $a=\|w^{t}-w^{*}\|$, $b=\|\nabla f(w^{t})\|$, we obtain
\[
\|\nabla f(w^{t})\|^{2}\ge \frac{2}{\eta E}\bigl(f(w^{t})-f(w^{*})\bigr)-\frac{1}{\eta^{2}E^{2}}\|w^{t}-w^{*}\|^{2}. \tag{17}
\]

\medskip
\noindent\underline{[Step 14]} \textbf{Plug (17) into (16).}  
After rearranging terms and taking total expectation,
\[
\begin{aligned}
\E\!\bigl[f(w^{t+1})-f(w^{*})\bigr]
&\le \E\!\bigl[f(w^{t})-f(w^{*})\bigr]
   -\frac{\eta E}{2}\Bigl(\frac{2}{\eta E}\bigl(f(w^{t})-f(w^{*})\bigr)-\frac{1}{\eta^{2}E^{2}}\|w^{t}-w^{*}\|^{2}\Bigr)\\
&\qquad+ \frac{\eta E}{2}\Delta^{2}
   +2L\eta^{2}E^{2}\Delta^{2}+ \frac{L\eta^{2}E^{2}\sigma^{2}}{K}\\
&= \E\!\bigl[f(w^{t})-f(w^{*})\bigr]
   -\bigl(f(w^{t})-f(w^{*})\bigr)
   +\frac{1}{2\eta E}\|w^{t}-w^{*}\|^{2}\\
&\qquad+ \frac{\eta E}{2}\Delta^{2}
   +2L\eta^{2}E^{2}\Delta^{2}+ \frac{L\eta^{2}E^{2}\sigma^{2}}{K}.
\end{aligned} \tag{18}
\]

\medskip
\noindent\underline{[Step 15]} \textbf{Telescoping sum.}  
Rearrange (18) to isolate the distance term:
\[
\frac{1}{2\eta E}\|w^{t}-w^{*}\|^{2}
\ge \E\!\bigl[f(w^{t})-f(w^{*})\bigr]
   -\E\!\bigl[f(w^{t+1})-f(w^{*})\bigr]
   - C,
\]
where
\[
C:=\frac{\eta E}{2}\Delta^{2}+2L\eta^{2}E^{2}\Delta^{2}+ \frac{L\eta^{2}E^{2}\sigma^{2}}{K}.
\]
Summing over $t=0,\dots,T-1$ yields
\[
\frac{1}{2\eta E}\|w^{0}-w^{*}\|^{2}
\ge \sum_{t=0}^{T-1}\E\!\bigl[f(w^{t})-f(w^{*})\bigr] - \underbrace{\sum_{t=0}^{T-1}\E\!\bigl[f(w^{t+1})-f(w^{*})\bigr]}_{\text{telescopes to } \E[f(w^{T})-f(w^{*})]\ge0}
   - T C .
\]
Since the last term is non‑negative, we drop it and obtain
\[
\sum_{t=0}^{T-1}\E\!\bigl[f(w^{t})-f(w^{*})\bigr]
\;\le\; \frac{\|w^{0}-w^{*}\|^{2}}{2\eta E}+ T C . \tag{19}
\]

\medskip
\noindent\underline{[Step 16]} \textbf{Divide by $T$ and simplify $C$.}  
Recall $C$ and use $E\ge1$:
\[
\begin{aligned}
\frac{1}{T}\sum_{t=0}^{T-1}\E\!\bigl[f(w^{t})-f(w^{*})\bigr]
&\le \frac{\|w^{0}-w^{*}\|^{2}}{2\eta E\,T}
   +\frac{\eta E}{2}\Delta^{2}+2L\eta^{2}E^{2}\Delta^{2}+ \frac{L\eta^{2}E^{2}\sigma^{2}}{K}\\
&\le \frac{\|w^{0}-w^{*}\|^{2}}{2\eta T}
   + L\eta\sigma^{2}\frac{1}{K}
   + L\eta\Delta^{2}E,
\end{aligned}
\]
where we used $\eta\le \frac{1}{2L(E+1)}$ to absorb the higher‑order term $2L\eta^{2}E^{2}\Delta^{2}$ into $L\eta\Delta^{2}E$.

Thus we have exactly the bound (4) claimed in Theorem \ref{thm:conv}. \qed

\section*{Rate Derivation (With Constants)}
From (4) set $\eta = \frac{c}{\sqrt{T}}$ with $c\le \frac{1}{2L(E+1)}$.
Then
\[
\frac{\|w^{0}-w^{*}\|^{2}}{2\eta T}= \frac{\|w^{0}-w^{*}\|^{2}}{2c}\,\frac{1}{\sqrt{T}},\qquad
L\eta\sigma^{2}/K = \frac{Lc\sigma^{2}}{K}\,\frac{1}{\sqrt{T}},\qquad
L\eta\Delta^{2}E = LcE\Delta^{2}\,\frac{1}{\sqrt{T}}.
\]
Hence
\[
\frac{1}{T}\sum_{t=0}^{T-1}\E\!\bigl[f(w^{t})-f(w^{*})\bigr]
= O\!\Bigl(\frac{1}{\sqrt{T}}\bigl(1+\tfrac{E\Delta^{2}}{K}+\tfrac{\sigma^{2}}{K}\bigr)\Bigr).
\]

\section*{Special Cases}
\begin{itemize}[leftmargin=*]
  \item \textbf{Homogeneous data ($\Delta=0$).} The drift term vanishes and we recover the classic $O(1/\sqrt{T})$ rate of centralized SGD with variance reduced by a factor $K$.
  \item \textbf{Single local step ($E=1$).} The bound becomes identical to standard distributed SGD, confirming that FedAvg with $E=1$ does not suffer extra drift.
  \item \textbf{Large $E$ with identical data.} When $\Delta=0$, increasing $E$ does not affect the asymptotic rate; the algorithm enjoys communication efficiency without loss of statistical accuracy.
\end{itemize}

\begin{thebibliography}{9}
\bibitem{Nesterov2004}
Y.~Nesterov, \emph{Introductory Lectures on Convex Optimization: A Basic Course}, Kluwer Academic Publishers, 2004.
\end{thebibliography}

\end{document}